\documentclass[UTF8,a4paper,12pt]{ctexart}

\usepackage{geometry}
\usepackage{graphicx}
\usepackage{booktabs}
\usepackage{amsmath}
\usepackage{hyperref}
\usepackage{listings}
\usepackage{xcolor}

\geometry{left=2.5cm,right=2.5cm,top=2.5cm,bottom=2.5cm}
\hypersetup{colorlinks=true,linkcolor=blue,citecolor=blue,urlcolor=blue}

\lstset{
  basicstyle=\ttfamily\small,
  breaklines=true,
  frame=single,
  columns=fullflexible
}

\title{基于 VoxCPM 的可控音色语音合成与语音聊天系统实现}
\author{姓名：\underline{\hspace{4cm}} \quad 学号：\underline{\hspace{4cm}}}
\date{\today}

\begin{document}
\maketitle

\begin{abstract}
本文面向语音信号处理课程大作业，围绕开源神经语音合成模型 VoxCPM/VoxCPM2 设计并实现了一个可控音色语音合成与语音聊天系统。项目重点在本地语音模型的推理与适配：通过自然语言音色描述、参考音频克隆和 LoRA 说话人微调，实现从文本条件到目标语音波形的生成；对话模型仅用于产生待朗读文本和语气条件，不参与语音信号建模。系统实现包括音频数据预处理、VoxCPM 条件输入封装、推理参数控制、参考音频复用、训练 manifest 构建和 LoRA 配置生成。论文从连续语音表征、层次化语义-声学建模、音色控制与实验评价四个方面分析系统。实验部分计划从实时率、可懂度、音色相似度、主观自然度和不同推理步数下的速度质量权衡进行评估。硬件分析表明，4GB 显存笔记本不适合舒适运行 VoxCPM2，需要采用 CPU 推理、低资源 VoxCPM-0.5B 或更高显存设备。

\textbf{关键词：} 语音合成；VoxCPM；声音克隆；音色控制；LoRA 微调
\end{abstract}

\section{引言}
随着神经网络声学模型、扩散模型和大规模语音生成模型的发展，语音合成已经从传统拼接合成、统计参数合成发展到能够进行零样本声音克隆和自然语言风格控制的生成式系统。课程项目关注深度学习与语音信号处理的结合，本文选择“可控音色语音合成”作为核心任务，并将其放入语音聊天场景中展示：系统接收待朗读文本和语气控制条件，再由本地 VoxCPM 模型生成目标音色的语音波形。

本项目的目标不是调用在线语音合成服务，而是复现并工程化一个开源本地语音模型。系统保留 VoxCPM 上游模型作为核心声学生成后端，在项目侧实现音色控制参数封装、参考音频管理、数据预处理、LoRA 配置生成和命令行推理入口。聊天文本可以由任意对话模型生成，但论文的技术重点始终是本地语音生成链路：如何组织文本、语气、参考音频和微调权重，使 VoxCPM 输出稳定、可控且可比较的语音。

\section{国内外研究现状}
早期语音合成常使用拼接合成和统计参数合成，近年来主流方法转向基于深度神经网络的端到端 TTS。Tacotron、FastSpeech、VITS 等模型推动了从文本到梅尔谱再到声码器的神经合成流程。近年的大规模语音生成模型进一步引入离散语音 token、扩散模型、流匹配和语言模型骨干，使系统具备零样本声音克隆、多说话人控制和跨语言生成能力。

VoxCPM 的核心特点是 tokenizer-free，即不依赖外部离散语音分词器，而是在连续语音表征空间中通过层次化语义-声学建模生成语音 \cite{zhou2025voxcpm}。VoxCPM2 在此基础上扩展到 2B 参数、多语言、自然语言音色设计、可控声音克隆和 48kHz 输出 \cite{voxcpm2_2026}。其官方文档给出了 LocEnc、TSLM、RALM、LocDiT 的四阶段结构，这为课程项目提供了一个具备研究价值的开源语音模型基础。

在交互系统中，对话模型只提供待朗读文本和语气描述，可由 DeepSeek 或其他兼容接口生成 \cite{deepseekai2025r1}。该部分属于上层文本条件生成，不改变语音模型的训练与推理机制。因此本文后续重点讨论 VoxCPM 的本地推理、音色控制、参考音频条件和 LoRA 适配，而不是对话 API 本身。

\section{方法}
\subsection{系统总体结构}
系统由四个模块组成：文本与语气条件生成、VoxCPM 本地语音合成、音色/参考音频控制、实验与微调工具链。给定待朗读文本 $y$、语气提示词 $p$、基础音色描述 $s$ 和可选参考音频 $a_r$，系统首先构造 VoxCPM 的条件文本：
\[
\tilde{y} = \mathrm{concat}((s, p), y).
\]
其中括号内的自然语言描述用于控制音色、情绪、语速和停顿。随后 VoxCPM 根据文本条件、可选参考音频和可选提示音频转写生成波形：
\[
\hat{a} = g_{\phi}(\tilde{y}, a_r, a_p, t_p).
\]
其中 $a_p$ 和 $t_p$ 分别表示极致克隆模式下的提示音频与对应转写文本。项目侧代码不修改 VoxCPM 主体结构，而是实现条件组织、模型加载、输出保存和实验参数控制。

\begin{figure}[htbp]
  \centering
  \fbox{\parbox{0.88\linewidth}{\centering 图 1 占位：系统流程图。请使用 paper/figures/prompts.md 中的 Prompt 1 生成后替换。}}
  \caption{基于本地 VoxCPM 的可控音色语音合成系统流程。}
  \label{fig:pipeline}
\end{figure}

\subsection{VoxCPM 连续表征语音生成}
VoxCPM 的关键思想是避免依赖外部离散语音分词器，直接在连续音频隐空间中建模语音。其整体流程可概括为四个阶段：局部编码器将音频片段映射到紧凑连续表征；文本-语义语言模型负责根据文本捕捉语义、韵律和停顿计划；残差声学语言模型补充细粒度声学信息；局部扩散 Transformer 在条件流匹配或扩散解码框架下生成连续音频 latent，最后由 AudioVAE 解码为波形。VoxCPM2 进一步支持 48kHz 输出、多语言合成、自然语言音色设计和参考音频可控克隆。

本项目的推理封装集中在三个参数：CFG 引导系数、扩散推理步数和最大生成长度。CFG 值影响条件约束强度，步数影响质量与速度，最大长度控制生成音频时长和显存占用。为了适应课程实验和低显存机器，默认关闭 denoiser 与编译优化，优先保证模型可加载和流程可复现。

\subsection{VoxCPM 语音合成与音色控制}
VoxCPM2 支持两种主要控制方式。第一，自然语言音色设计，即在合成文本前加入括号形式的描述，例如“年轻女性，声音温柔自然，语速中等”。第二，参考音频克隆，即提供参考语音，使模型生成接近参考说话人的音色。项目将用户指定的基础音色或参考音频作为全局设置，在连续对话中保持不变；每轮语气提示词只调整情绪、语速、停顿和重音，不改变说话人身份。

项目在 \texttt{VoxCPMSynthesizer} 中统一封装模型加载、LoRA 权重加载、条件文本拼接和音频文件输出。对单句合成，输入为文本和音色描述；对参考音频克隆，输入额外包含 \texttt{reference\_audio}；对极致克隆，输入包含 \texttt{prompt\_audio} 与 \texttt{prompt\_text}。最终输出为 WAV 文件，采样率由 VoxCPM 模型本身决定。

\subsection{LoRA 微调}
为了适配特定说话人，项目提供 CSV/JSONL 到 VoxCPM manifest 的预处理脚本，并生成上游 LoRA 微调配置。预处理阶段将原始音频统一转换为单声道 16kHz WAV，过滤过短或过长样本，并写出包含音频路径、转写文本、时长和可选说话人编号的 JSONL manifest。LoRA 微调只更新低秩适配参数，较适合课程实验中的小规模自录数据。微调完成后，推理阶段通过 \texttt{lora\_path} 加载适配权重，从而在保留基础模型能力的同时增强目标说话人音色。

\section{实验设置与结果分析}
\subsection{实验环境}
当前开发机器为 Intel i7-12650H、13GB RAM、NVIDIA RTX 3050 Laptop GPU 4GB。软件环境为 Python 3.11、uv、PyTorch 2.5+、Transformers 和 VoxCPM 上游代码。

\subsection{数据与任务}
实验计划使用两类数据。第一类为自录普通话语音，时长 5--10 分钟，按句切分并人工转写，用于 LoRA 适配。第二类为开源中文 TTS 数据集中的少量样本，用于对比不同音色控制设置下的合成效果。若使用公开数据集，需在最终提交中补充数据集许可证与下载地址。

\subsection{评价指标}
实验将记录以下指标：

\begin{itemize}
  \item 实时率 RTF：合成耗时除以生成语音时长，越低越好。
  \item 字错率/词错率：使用 ASR 转写合成音频并与目标文本比较。
  \item 音色相似度：使用说话人嵌入模型计算参考音频与生成音频相似度，或进行主观 MOS 打分。
  \item 主观自然度：从清晰度、自然度、情感一致性三个维度进行 1--5 分打分。
\end{itemize}

\subsection{初步结果表}
由于模型权重和实验音频尚未全部生成，表 \ref{tab:placeholder} 先给出最终论文需要填写的结果格式。跑通实验后应替换为真实数值。

\begin{table}[htbp]
  \centering
  \caption{不同设置下的语音生成结果占位表}
  \label{tab:placeholder}
  \begin{tabular}{lllll}
    \toprule
    模型设置 & 推理设备 & RTF $\downarrow$ & 主观自然度 $\uparrow$ & 备注 \\
    \midrule
    VoxCPM2 + 音色设计 & CPU & 待测 & 待测 & 质量优先 \\
    VoxCPM-0.5B + Prompt 克隆 & CPU/GPU & 待测 & 待测 & 低资源演示 \\
    VoxCPM2 + LoRA & CPU/高显存 GPU & 待测 & 待测 & 个性化音色 \\
    \bottomrule
  \end{tabular}
\end{table}

\section{结论}
本文实现了一个基于本地 VoxCPM 的可控音色语音合成与语音聊天系统，覆盖了语音合成推理、音色控制、参考音频克隆、数据预处理和 LoRA 微调配置。系统将对话模型限制为文本和语气条件来源，语音生成模型在本地运行，符合课程对深度学习与语音信号处理结合的要求。硬件分析表明，4GB 显存不足以舒适运行 VoxCPM2，后续实验需要在低资源模型、CPU 推理和高显存设备之间权衡。下一步工作是补充真实音频样例、定量 RTF/ASR 指标、音色相似度、主观评分表，并将生成的科研图替换到论文中。

\bibliographystyle{plain}
\bibliography{references}

\appendix
\section{核心运行命令}
\begin{lstlisting}[language=bash]
uv sync --extra dev
uv pip install -e VoxCPM
export DEEPSEEK_API_KEY="sk-your-key-here"
uv run voiceagent synth --tts-model models/VoxCPM2 --tts-device cpu \
  --voice "young female voice, warm and natural tone, medium pace" \
  --text "你好，这是本地语音聊天系统。" --output outputs/demo.wav
\end{lstlisting}

\end{document}
